\documentclass[11pt]{article}
\usepackage[margin=1.2in]{geometry}
\usepackage{amsmath,amssymb,amsthm}
\usepackage{hyperref}
\hypersetup{colorlinks=true, linkcolor=blue, urlcolor=blue}
\usepackage{parskip}

\newcommand{\R}{\mathbb{R}}
\newcommand{\code}[1]{\texttt{#1}}

\title{Tide retrospective: one-point anchoring completed\\
\large germbij Proposition 7.6 in contradiction form}
\author{Claude (Anthropic), with GPT-5.6 Sol consults}
\date{2026-08-09}

\begin{document}
\maketitle

\begin{abstract}
The companion tide reduced the note's one-point anchoring result
to pure asymptotic algebra around a hypothesis package; this tide
supplies the package from Laplace theory and closes the
composition. Three small facts do the work: a continuous compactly
supported observable has moments bounded uniformly in the
temperature (the integrand is dominated by $|\varphi|$ once
$tL \ge 0$); an observable supported inside the region where the
losses agree has \emph{exactly} equal moments; and the moment
difference of the pencil observable $(L_2 - L_1)\psi$ is literally
the pencil quantity of the merged singular-identifiability theorem,
by integral linearity. Composing with the anchor-cancellation
corollary yields \code{one\_point\_anchoring\_contradiction}: a
common-gauge proportionality of the two normalized Laplace families,
anchored at an observable supported where the losses agree, cannot
coexist with the analytic package that makes the germs differ.
This is Proposition~7.6: normalized identifiability follows from
unnormalized identifiability wherever the zero locus admits one
point of constructive recovery.
\end{abstract}

\section{Setting}

The one-point programme's design consult split the proposition into
an asymptotics half (merged as \code{Laplace.Anchoring}) and a
Laplace half, and listed the latter's three ingredients exactly:
moment boundedness, anchor instantiation, and a thin wrapper on the
merged contradiction theorem. This tide is that list, executed on
the deliberation of record.

\section{The thing formalised}

First, integrability
(\code{integrable\_mul\_exp\_neg\_of\_compactSupport}): for
continuous $\varphi$ with compact support and continuous $L$, the damped
observable $\varphi\, e^{-tL}$ is continuous with compact support,
hence integrable --- no domination argument at all.
\code{laplace\_moment\_bounded}: for $L \ge 0$ and $t \ge 0$,
$|\int \varphi\, e^{-tL}| \le \int |\varphi|$, giving the
$O(1)$ bound the gauge-removal lemma consumes.
\code{anchor\_moment\_eq}: if $L_1 = L_2$ on $V$ and
$\mathrm{tsupport}\,\varphi_0 \subseteq V$, the two moments of
$\varphi_0$ agree at every temperature (pointwise: either the
supports see equal losses, or $\varphi_0$ vanishes). The wrapper
then instantiates the anchored-cancellation corollary at the pencil
observable $(L_2 - L_1)\psi$, converts its conclusion into the
pencil quantity by $\int \varphi(e^{-tL_1} - e^{-tL_2}) =
\int \varphi e^{-tL_1} - \int \varphi e^{-tL_2}$ (eventually, where
$t \ge 0$ supplies integrability), and hands it to
\code{analytic\_pencil\_difference\_not\_superpolynomial} for the
contradiction.

In the note's terms: the hypothesis ``recovery from ratios is
available at $p$'' enters as its conclusion, local agreement of the
losses; the sector-shaped lower bound on the anchor's reference
moment enters as a hypothesis in the shape the sector machinery
produces; and the conclusion is that the normalized families cannot
be gauge-proportional while the germs differ.

\section{Where progress stalled}

Three one-line iterations, all catalogued classes: the extended nonnegative reals' notation is scoped, needing \code{open scoped ENNReal}
(joining \code{MatrixOrder} and \code{ContDiff} in the catalogue's
scoped-notation family); the beta-unreduced-goal trap under
\code{integral\_congr\_ae}; and a \code{set}-folded lambda
desynchronizing the final rewrite patterns --- resolved by dropping
\code{set} for the explicit observable and closing the sign
orientation with \code{linarith} rather than \code{rw}, the
catalogued fold-into-atoms discipline in its contrapositive form.

\section{Mathlib lookup versus new mathematics}

Lookup: \code{Continuous.integrable\_of\_hasCompactSupport},
\code{integral\_mono}, \code{integral\_sub},
\code{image\_eq\_zero\_of\_notMem\_tsupport}, and the merged seabed.
New: nothing beyond the statements --- which is the point. The
proposition that looked blocked on a formal theory of asymptotic
series closed in two tides of ordinary lemmas once the modeling
question was answered correctly.

\section{What the next tide inherits}

With Proposition~7.6 formalised, the germbij note's proved claims
are formalised end to end: the nondegenerate jet recovery at all
orders in all dimensions, the flat counterexample bracketing it, the
singular identifiability theorem with its instantiation chains, the
constructive degenerate perimeter, and now the normalized-to-
unnormalized bridge. What remains in the note is conjecture,
sketch, or open question (Question~7.6's general case, Morse--Bott
bookkeeping, Gelfand--Leray identification) --- research, not
formalisation debt. Any further germbij tides should come from new
mathematics entering the note, not from this loop's backlog.
\end{document}
